% META: {"id":"1SPE-SECDEG-CO-103","chapitre":"1SPE-SECOND-DEGRE","type_objet":"corrige","exercice_id":"1SPE-SECDEG-EX-103","status":"generated"}
% BEGIN-VERIFY
% from math import sqrt
% delta = 2 * 2 - 4 * 1 * (-0.21)
% assert abs(delta - 4.84) < 1e-12 and abs(sqrt(delta) - 2.2) < 1e-12
% t1 = (-2 - 2.2) / 2
% t2 = (-2 + 2.2) / 2
% assert abs(t1 + 2.1) < 1e-12 and abs(t2 - 0.1) < 1e-12
% assert abs((1 + t2) ** 2 - 1.21) < 1e-12
% assert 1 + t1 == -1.1 and 1 + t1 < 0
% assert abs(1.3 * 0.8 - 1.04) < 1e-12
% equivalent = sqrt(1.04) - 1
% assert round(equivalent, 4) == 0.0198
% assert abs((1 + equivalent) ** 2 - 1.04) < 1e-12
% END-VERIFY
\begin{corrige}{1SPE-SECDEG-EX-103}
\begin{enumerate}
\item Deux évolutions successives de taux $t$ multiplient la population par
  $(1+t)$ deux fois, donc par $(1+t)^2$. Une hausse globale de $21\,\%$
  correspond au coefficient $1{,}21$, d'où $(1+t)^2 = 1{,}21$.

\item En développant : $1 + 2t + t^2 = 1{,}21$, soit
  $t^2 + 2t - 0{,}21 = 0$.
  \[
    \Delta = 2^2 - 4\times1\times(-0{,}21) = 4 + 0{,}84 = 4{,}84,
    \qquad \sqrt{\Delta} = 2{,}2 .
  \]
  Les racines sont
  $t_1 = \dfrac{-2 - 2{,}2}{2} = -2{,}1$ et
  $t_2 = \dfrac{-2 + 2{,}2}{2} = 0{,}1$.

\item On écarte $t_1 = -2{,}1$ : le coefficient multiplicateur associé vaut
  $1 + t_1 = -1{,}1$, or un coefficient multiplicateur est toujours
  \textbf{positif}. Le taux annuel est donc $t_2 = 0{,}1$, soit $+10\,\%$
  par an. Vérification : $1{,}1^2 = 1{,}21$. \checkmark

\item Les coefficients se multiplient :
  $1{,}30 \times 0{,}80 = 1{,}04$, soit une hausse globale de $4\,\%$.
  Le taux annuel équivalent $t$ vérifie $(1+t)^2 = 1{,}04$, d'où
  $1 + t = \sqrt{1{,}04}$ et
  \[
    t = \sqrt{1{,}04} - 1 \approx 0{,}0198,
  \]
  soit environ $+1{,}98\,\%$ par an -- et non la moyenne
  $\dfrac{30 - 20}{2} = 5\,\%$.
\end{enumerate}
\end{corrige}
